\documentclass[11pt]{article}
\usepackage[utf8]{inputenc}
\usepackage{amsmath,amssymb}
\usepackage{hyperref}
\title{Efficient Landslide Susceptibility Mapping: A Resource-Aware Approach}
\author{Anonymous}
\date{}
\begin{document}
\maketitle

\begin{abstract}
This paper studies Landslide Susceptibility Mapping. Grounded in a real, citable literature \cite{ref1,ref2,ref3,ref4,ref5,ref6,ref7,ref8},
we propose a focused method and report \textbf{illustrative} experiments.
All references below are real and verifiable (OpenAlex/arXiv); the reported
numbers are simulated to illustrate intended behaviour.
\end{abstract}

\section{Introduction}
Landslide Susceptibility Mapping is an active research area. This work targets a concrete gap identified from
the recent literature and contributes a method together with an evaluation protocol.

\section{Related Work (real literature)}
The following works, retrieved from public bibliographic APIs, frame our study:
\begin{itemize}
\item Guzzetti et al. (2012) — "Landslide inventory maps: New tools for an old problem", Earth-Science Reviews (cited 2089).
\item Ayalew et al. (2004) — "The application of GIS-based logistic regression for landslide susceptibility mapping in the Kakuda-Yahiko Mountains, Central Japan", Geomorphology (cited 1921).
\item Fell et al. (2008) — "Guidelines for landslide susceptibility, hazard and risk zoning for land use planning", Engineering Geology (cited 1476).
\item Bui et al. (2015) — "Spatial prediction models for shallow landslide hazards: a comparative assessment of the efficacy of support vector machines, artificial neural networks, kernel logistic regression, and logistic model tree", Landslides (cited 1260).
\item Pradhan (2012) — "A comparative study on the predictive ability of the decision tree, support vector machine and neuro-fuzzy models in landslide susceptibility mapping using GIS", Computers \& Geosciences (cited 1209).
\item Merghadi et al. (2020) — "Machine learning methods for landslide susceptibility studies: A comparative overview of algorithm performance", Earth-Science Reviews (cited 1163).
\end{itemize}

\section{Method}
We reduce the dominant computational cost via adaptive resource allocation, activating only the components needed per input. We instantiate this idea for Landslide Susceptibility Mapping and analyse its cost/quality trade-off.

\section{Experiments (Illustrative)}
\begin{center}
\begin{tabular}{lcc}
System & Primary Metric & Efficiency \\ \hline
Strong baseline & 0.812 & 1.00$\times$ \\
Ablation & 0.828 & 0.92$\times$ \\
\textbf{Ours} & \textbf{0.851} & \textbf{0.71$\times$}
\end{tabular}
\end{center}
\emph{Metrics simulated to illustrate the intended effect; not measured.}

\section{Conclusion}
We connected a real literature to a concrete method for Landslide Susceptibility Mapping. Future work will
validate the illustrative results with real experiments.

\begin{thebibliography}{9}
\bibitem{ref1} Fausto Guzzetti, Alessandro Mondini, Mauro Cardinali et al.. Landslide inventory maps: New tools for an old problem. \emph{Earth-Science Reviews}, 2012. DOI:10.1016/j.earscirev.2012.02.001
\bibitem{ref2} Lulseged Ayalew, Hiromitsu Yamagishi. The application of GIS-based logistic regression for landslide susceptibility mapping in the Kakuda-Yahiko Mountains, Central Japan. \emph{Geomorphology}, 2004. DOI:10.1016/j.geomorph.2004.06.010
\bibitem{ref3} Robin Fell, Jordi Corominas, C. Bonnard et al.. Guidelines for landslide susceptibility, hazard and risk zoning for land use planning. \emph{Engineering Geology}, 2008. DOI:10.1016/j.enggeo.2008.03.022
\bibitem{ref4} Dieu Tien Bui, Trần Anh Tuấn, Harald Klempe et al.. Spatial prediction models for shallow landslide hazards: a comparative assessment of the efficacy of support vector machines, artificial neural networks, kernel logistic regression, and logistic model tree. \emph{Landslides}, 2015. DOI:10.1007/s10346-015-0557-6
\bibitem{ref5} Biswajeet Pradhan. A comparative study on the predictive ability of the decision tree, support vector machine and neuro-fuzzy models in landslide susceptibility mapping using GIS. \emph{Computers \& Geosciences}, 2012. DOI:10.1016/j.cageo.2012.08.023
\bibitem{ref6} Abdelaziz Merghadi, Ali P. Yunus, Jie Dou et al.. Machine learning methods for landslide susceptibility studies: A comparative overview of algorithm performance. \emph{Earth-Science Reviews}, 2020. DOI:10.1016/j.earscirev.2020.103225
\bibitem{ref7} Michel Jaboyedoff, Thierry Oppikofer, Antonio Abellán et al.. Use of LIDAR in landslide investigations: a review. \emph{Natural Hazards}, 2010. DOI:10.1007/s11069-010-9634-2
\bibitem{ref8} Hamid Reza Pourghasemi, Biswajeet Pradhan, Candan Gökçeoğlu. Application of fuzzy logic and analytical hierarchy process (AHP) to landslide susceptibility mapping at Haraz watershed, Iran. \emph{Natural Hazards}, 2012. DOI:10.1007/s11069-012-0217-2
\end{thebibliography}
\end{document}
